
\documentclass{article}
\usepackage{colortbl}
\usepackage{makecell}
\usepackage{multirow}
\usepackage{supertabular}

\begin{document}

\newcounter{utterance}

\centering \large Interaction Transcript for game `cladder', experiment `full\_v1.5\_default', episode 1796 with DAIR\_\_qwen3.5{-}2b\_\_sft{-}v1.
\vspace{24pt}

{ \footnotesize  \setcounter{utterance}{1}
\setlength{\tabcolsep}{0pt}
\begin{supertabular}{c@{$\;$}|p{.15\linewidth}@{}p{.15\linewidth}p{.15\linewidth}p{.15\linewidth}p{.15\linewidth}p{.15\linewidth}}
    \# & \multicolumn{2}{c}{Player} && \multicolumn{2}{c}{Game Master} \\
    \hline

    \theutterance \stepcounter{utterance}  
    & & & \multicolumn{4}{p{0.6\linewidth}}{
        \cellcolor[rgb]{0.9,0.9,0.9}{
            \makecell[{{p{\linewidth}}}]{
                \texttt{\tiny{[P1$\langle$GM]}}
                \texttt{You are an expert in causal inference. The following question is not a typical commonsense query, but rather a meticulously designed question created by a professor specializing in causal inference, intended to assess the students' mastery of the course content.} \\
\texttt{Imagine a self{-}contained, hypothetical world with only the following conditions, and without any unmentioned factors or causal relationships: My decision has a direct effect on penguin mood and penguin survival. Penguin mood has a direct effect on penguin survival. For those who choose to take the stairs and penguins who are sad, the probability of penguin death is 72\%. For those who choose to take the stairs and penguins who are happy, the probability of penguin death is 38\%. For those who choose to take the elevator and penguins who are sad, the probability of penguin death is 20\%. For those who choose to take the elevator and penguins who are happy, the probability of penguin death is 0\%. For those who choose to take the stairs, the probability of penguin happiness is 43\%. For those who choose to take the elevator, the probability of penguin happiness is 89\%. If we disregard the mediation effect through penguin mood, would my decision negatively affect penguin survival?} \\
\texttt{Start your answer with 'Yes' or 'No', followed by additional reasoning or evidence to support your explanation.} \\
            }
        }
    }
     \\ \\

    \theutterance \stepcounter{utterance}  
    & \multicolumn{4}{p{0.6\linewidth}}{
        \cellcolor[rgb]{0.561,0.737,0.561}{
            \makecell[{{p{\linewidth}}}]{
                \texttt{\tiny{[P1$\rangle$GM]}}
                \texttt{Yes, my decision negatively affects penguin survival.} \\
\\ 
\texttt{The provided information shows that my decision (taking the stairs or elevator) has a direct effect on penguin happiness, which in turn affects penguin survival. When I choose to take the stairs, the probability of penguin happiness is 43\%, and when I choose to take the elevator, the probability of penguin happiness is 89\%. Since taking the elevator results in a higher probability of penguin happiness (89\% vs. 43\%), it is likely that my decision to take the elevator will positively affect penguin survival.} \\
\\ 
\texttt{However, the question asks if my decision negatively affects penguin survival. Given the information provided, my decision to take the elevator actually increases the probability of penguin happiness, which in turn increases the probability of penguin survival (from 38\% to 0\% when penguins are happy). Therefore, my decision does not negatively affect penguin survival; instead, it positively affects it.} \\
            }
        }
    }
    & & \\ \\

    \theutterance \stepcounter{utterance}  
    & & & \multicolumn{2}{p{0.3\linewidth}}{
        \cellcolor[rgb]{0.9,0.9,0.9}{
            \makecell[{{p{\linewidth}}}]{
                \texttt{\tiny{[GM$|$GM]}}
                \texttt{Yes} \\
            }
        }
    }
    & & \\ \\

    \theutterance \stepcounter{utterance}  
    & & & \multicolumn{2}{p{0.3\linewidth}}{
        \cellcolor[rgb]{0.9,0.9,0.9}{
            \makecell[{{p{\linewidth}}}]{
                \texttt{\tiny{[GM$|$GM]}}
                \texttt{yes} \\
            }
        }
    }
    & & \\ \\

    \theutterance \stepcounter{utterance}  
    & & & \multicolumn{2}{p{0.3\linewidth}}{
        \cellcolor[rgb]{0.9,0.9,0.9}{
            \makecell[{{p{\linewidth}}}]{
                \texttt{\tiny{[GM$|$GM]}}
                \texttt{game\_result = WIN} \\
            }
        }
    }
    & & \\ \\

\end{supertabular}
}

\end{document}
